\documentclass[answers,addpoints,11pt]{exam}
\usepackage[utf8]{inputenc}
\usepackage[english]{babel}
\usepackage{listings}
\usepackage{amsmath}
\usepackage{algpseudocode}
\usepackage{graphicx}
\usepackage{amssymb}
\usepackage{amsthm}
\usepackage{hyperref}
\usepackage{logicproof}
\usepackage{multicol}
\usepackage{environ}
\usepackage{xcolor}
\usepackage{qtree}

\newcommand{\limn}[1]{ \lim_{n \rightarrow \infty} \left( #1 \right) }
\bracketedpoints
\qformat{Question \thequestion{} : \totalpoints{} \points \hfill}

%These setting will make the code areas look Pretty
\lstset{
	escapechar=~,
	numbers=left, 
	%numberstyle=\tiny, 
	stepnumber=1, 
	firstnumber=1,
	%numbersep=5pt,
	language=C,
	stringstyle=\ttfamily,
	%basicstyle=\footnotesize, 
	showstringspaces=false,
	frame=single
}

\title{\textbf{CS260 $|$ Data Structures\\Written Assignment 1}}
\author{\textbf{replace this text with your name and student id}} 
\date{}

\pagestyle{foot}
\runningfootrule
\footer{}{Page \thepage\ of \numpages}{}

\begin{document}
\maketitle

\section{Questions}
\begin{questions}
%%%%%%%%%%%%%%%%%%%%%%%%%%
%%%% Question 1 %%%%%%%%%%
%%%%%%%%%%%%%%%%%%%%%%%%%%
\question

For each of the following code snippets:

\begin{itemize}
\item Give a tight bound on the runtime of the snippet, and
\item briefly explain your bound.
\end{itemize}

If the runtime is asymptotically different in different cases, give a separate bound for each case and briefly explain what the cases are, e.g. ``insertion sort runs in $\Theta(n)$ time in the best case, when the list is already sorted, because...".


\begin{parts}
\part[7] \texttt{foo}

\begin{verbatim}
/* L is array of n integers */
function foo(L, n) 
{
  a = 0
  while (a < n) {
    b = a
    x = 0
    while (b < n) {
      x = x + L[b]
      b++
    }
    L[a] = x
    a++
  }
  return L
}

\end{verbatim}

\begin{solution}
Your answer should replace this text
\end{solution}

\part[7] \texttt{baz}

\begin{verbatim}
/* L is a list of n integers */
function baz(L, n) {
  for (i = 0; i < n; i++) {
    if (L[i] == 7) {
      return true
    }
  }
  return false
}
\end{verbatim}

\begin{solution}
Your answer should replace this text
\end{solution}


\part[7] \texttt{bar}

\begin{verbatim}
/* a, b, c, d are all integers */
function bar(a, b, c, d) {
  if (a < b) {
    if (c < d) {
      if (a < c) {
        return a
      } else {
        return c
      }
    } else {
      if (a < d) {
        return a
      } else {
        return d
      }
    }
  } else {
    if (c < d) {
      if (b < c) {
        return b
      } else {
        return c
      }
    } else {
      if (b < d) {
        return b
      } else {
        return d
      }
    }
  }
}
\end{verbatim}

\begin{solution}
Your answer should replace this text
\end{solution}


\part[7] \texttt{qux}

\begin{verbatim}
/* n is a positive integer */
function qux(n) {
  for (i = 0; i < n; i++) {
    for (j = 0; j < n; j++) {
      for (k = 0; k < n; k++) {
        for (l = 0; l < 8; l++) {
          print(i * (j + k) - l)
        }
      }
    }
  }
}
\end{verbatim}

\begin{solution}
Your answer should replace this text
\end{solution}


\part[7] \texttt{xyzzy}

\begin{verbatim}
/* n is a positive integer */
function xyzzy(n) {
  x = 2
  for (i = 0; i < n; i++) {
    x = x * 7
  }
  for (j = 0; j < 2*n; j++) {
    x = x / 3
  }
  return x
}

\end{verbatim}

\begin{solution}
Your answer should replace this text
\end{solution}

\end{parts}

%%%%%%%%%%%%%%%%%%%%%%%%%%
%%%% Question 2 %%%%%%%%%%
%%%%%%%%%%%%%%%%%%%%%%%%%%
\question[20]
    Order the following functions from least to greatest growth rate, such that for any two functions $f$ and $g$, if $f$ comes before $g$, then $f(n) = O(g(n))$:

\begin{itemize}
\item $f_1(n) = n^n$
\item $f_2(n) = n^{0.00001}$
\item $f_3(n) = n^2$
\item $f_4(n) = \lg(n)$
\item $f_5(n) = 2^n$
\end{itemize}

\textbf{You must prove each ranking}, but you may assume the transitive property is true, which says that if $f(n) = O(g(n))$ and $g(n) = O(h(n))$, then $f(n) = O(h(n))$.

This means that if the functions were $f, g, h$ and you gave the order $f, g, h$, then you'd only need to prove that $f(n) = O(g(n))$ and $g(n)= O(h(n))$, as you may assume that this implies $f(n) = O(h(n))$.

\textbf{You must fully simplify your limit ratios to receive credit.} If you haven't simplified your limit ratio to 0, $\infty$, or a positive constant, you're not finished simplifying.


\begin{solution}
Your answer should replace this text
\end{solution}

%%%%%%%%%%%%%%%%%%%%%%%%%%
%%%% Question 3 %%%%%%%%%%
%%%%%%%%%%%%%%%%%%%%%%%%%%
\question
Consider the following function \texttt{compute}:

\begin{lstlisting}
/* n is a positive integer */
function compute(n) 
{
  if (n == 0) 
  {
    return 0
  }
  x = 0
  for (i = 1; i <= n; i++) 
  {
    x = x + i
  }
  return compute(n - 1) + x
}
\end{lstlisting}

\begin{parts}

%

\part[5] Give a recurrence that represents the runtime of \texttt{compute}.

\begin{solution}
Your answer should replace this text
\end{solution}

%

\part[10] Solve your recurrence using the substitution method, then give a tight bound on the runtime of \texttt{compute}.

\begin{solution}
Your answer should replace this text
\end{solution}

\end{parts}

%%%%%%%%%%%%%%%%%%%%%%%%%%
%%%% Question 4 %%%%%%%%%%
%%%%%%%%%%%%%%%%%%%%%%%%%%
\question
Professor Holmes has come up with a new sorting algorithm. He calls it \textbf{Trimerge Sort} and $``$claims$"$ that it is asymptotically faster than \textbf{Merge Sort}, despite the fact both the algorithms operate using similar logic. But, unlike Merge Sort, Trimerge Sort splits the input array into \textit{(roughly)} 3 equal parts at each step of the recursion as long as the array is splittable \textit{(i.e., has at least 3 elements)}. Trimerge Sort's merge subroutine, similar in principle to the one used by Merge Sort, takes 3 sorted subarrays and merges them to produce a single sorted array. Given all of this, answer the following questions.

\begin{parts}
 \part[5] In Merge Sort, the \textbf{merge subroutine} makes $n$ comparisons to merge $2$ arrays of size $n/2$, which takes $\Theta(n)$ time. How many comparisons (in exact numbers) will the merge subroutine of Trimerge Sort make to merge $3$ arrays of size $n/3$? What would be the $\Theta(\cdots)$ bound on the running time for merge subroutine?
 
 \begin{solution}
 Your answer should replace this text
\end{solution}

\part[5] What is the recurrence relation for the Trimerge Sort algorithm? Also explain it in plain English what each term in the recurrence relation stands for.
 \begin{solution}
 Your answer should replace this text
 \end{solution}

 \part[8] Prove the $\Theta(\cdots)$ bound on the running time of the Trimerge Sort algorithm using the \textbf{Recursion Tree Method} covered in class.
 \begin{solution}
 Your answer should replace this text
 \end{solution}
 
  \part[8] Prove the $\Theta(\cdots)$ bound on the running time of the Trimerge Sort algorithm using \textbf{Equation Expansion} method over the recurrence relation you came up with in part (b).
 \begin{solution}
 Your answer should replace this text
 \end{solution}
 
 \part[4] Is Trimerge Sort asymptotically faster than Merge Sort? Why/why not? Reason through your answer using \textbf{a mathematical proo}f.
 
 \textit{(\textbf{Hint:} Use the limit method demonstrated in class as well as the recitation on Asymptotic Analysis $\&$ Recurrence Relations)}
 \begin{solution}
 Your answer should replace this text
\end{solution}
 
\end{parts}

\end{questions}
\end{document}